\documentclass[10pt,letterpaper]{article}

\usepackage[T1]{fontenc}
\usepackage[margin=0.72in]{geometry}
\usepackage{amsmath,amssymb}
\usepackage{array}
\usepackage{booktabs}
\usepackage[dvipsnames]{xcolor}
\usepackage{enumitem}
\usepackage{hyperref}
\usepackage{listings}
\usepackage{longtable}
\usepackage{tabularx}

\hypersetup{
  colorlinks=true,
  linkcolor=MidnightBlue,
  urlcolor=MidnightBlue,
  pdftitle={TileGrad: north-star specification},
}

\lstset{
  language=Python,
  basicstyle=\ttfamily\small,
  columns=fullflexible,
  keepspaces=true,
  frame=single,
  framerule=0.2pt,
  xleftmargin=0.4em,
  xrightmargin=0.4em,
  aboveskip=0.5em,
  belowskip=0.5em,
}

\setlist[itemize]{leftmargin=1.5em,itemsep=1pt,topsep=3pt}
\setlist[enumerate]{leftmargin=1.7em,itemsep=1pt,topsep=3pt}
\newcolumntype{Y}{>{\raggedright\arraybackslash}X}
\newcommand{\TileGrad}{\textsc{TileGrad}}
\newcommand{\tinygrad}{\textsc{tinygrad}}
\newcommand{\op}[1]{\textsc{#1}}
\newcommand{\must}{\textbf{MUST}}
\newcommand{\mustnot}{\textbf{MUST NOT}}
\newcommand{\should}{\textbf{SHOULD}}
\newcommand{\may}{\textbf{MAY}}

\title{\TileGrad: a typed tile language over \tinygrad{} UOps\\
\large North-star specification, Version 1}
\author{TileGrad contributors}
\date{July 2026}

\begin{document}
\maketitle

\begin{abstract}
\TileGrad{} is a Python-embedded language for tiled compute kernels.
Programs declare tile shapes, memory placement, launch shape, and coarse
scheduling intent; the compiler owns lane-level distribution, synchronization,
and legal instruction selection.  \TileGrad{} lowers into \tinygrad{} UOps and
delegates backend code generation, compilation, and runtime execution to
\tinygrad{}.

This is a \emph{north-star contract}, not a snapshot of the current experimental
implementation.  It states what \TileGrad{} must become to be considered done,
which decisions are stable, and which surface is explicitly deferred.  Normative
text uses the keywords \must, \mustnot, \should, and \may.  Exact keyword
signatures, error-message wording, and numerical tolerances live in the API
reference and the conformance test suite respectively; this document defers to
them where it says so.
\end{abstract}

\tableofcontents

\section{Why \TileGrad{}}

\tinygrad{} compiles tensor graphs and runs them, but it has no language for an
\emph{author to write a custom kernel}.  Today the only way to add one is to
hand-roll bare UOps: lane indices, shared-memory barriers, swizzled layouts,
and WMMA operand scopes are all the user's problem.  Triton and TileLang solve
this in general, but they also own their own backends, schedulers, and
autotuners; they are large systems to bridge to \tinygrad{} from and large
systems to keep coherent.

\TileGrad{} takes a narrower position:

\begin{itemize}
  \item Make tiled kernels substantially easier to write than raw \tinygrad{}
        custom UOps, without inventing a new backend.
  \item Keep tile shape, memory placement, launch shape, and coarse scheduling
        intent visible in source and in debug IR.
  \item Provide one correctness path whose meaning is independent of target
        instruction selection; a scalar or SIMT fallback is always semantically
        equivalent to any selected native capability.
  \item Permit optimized lowerings without changing source-level semantics.
\end{itemize}

\TileGrad{} owns the typed tile IR, view and effect analysis, uniformity,
canonical logical layouts, synchronization planning, tile-op legalization, and
SIMT algorithms.  \tinygrad{} owns UOps, scalar and vector rendering,
device-specific code generation, compilation, launch, and outer tensor-graph
scheduling.  The compiler boundary between the two is the UOp interface
described in Section~\ref{sec:boundary}.

\section{Target workload}

The motivating workload is tiled GEMM and FlashAttention-style fused
attention.  If \TileGrad{} cannot express a clean FlashAttention forward pass,
its abstractions are wrong.  The canonical GEMM in Section~\ref{sec:program}
and the FlashAttention sketch in Section~\ref{sec:flash} are part of the
contract, not illustrations.

\subsection{Canonical GEMM}\label{sec:program}

The following spelling is the canonical V1 frontend example.  The frontend is a
Python JIT decorator; calling a decorated function specializes, compiles,
caches, and executes the kernel.  Eager or lazy evaluation of the body is not
inferred from the return type.

\begin{lstlisting}
@tilegrad.jit
def matmul(A, B, *,
           BM: T.Const[int] = 64,
           BN: T.Const[int] = 64,
           BK: T.Const[int] = 32):
  M, N, K = T.const("M", "N", "K")
  A: T.Tensor[(M, K), T.float16]
  B: T.Tensor[(K, N), T.float16]
  C = T.empty((M, N), T.float16)

  with T.Kernel(T.ceildiv(N, BN), T.ceildiv(M, BM),
               threads=128) as (bx, by):
    A_shared = T.shared((BM, BK), T.float16)
    B_shared = T.shared((BK, BN), T.float16)
    accum = T.fragment((BM, BN), T.float32)
    T.clear(accum)

    for ko in T.Pipelined(T.ceildiv(K, BK), stages=3):
      T.copy(A.tile((by*BM, ko*BK), (BM, BK)), A_shared, pad=0)
      T.copy(B.tile((ko*BK, bx*BN), (BK, BN)), B_shared, pad=0)
      T.mma(A_shared, B_shared, accum)

    T.copy(accum, C.tile((by*BM, bx*BN), (BM, BN)))
  return C
\end{lstlisting}

The reader who only reads this example should already understand the model:
runtime tensors are declared with shape and dtype; specialization symbols are
named once and bind to runtime dimensions; the launch region yields grid
coordinates; managed shared and collective fragment allocations live inside
that region; tile ops (\texttt{T.copy}, \texttt{T.clear}, \texttt{T.mma})
exchange typed tile regions; pipelining is a hint.

\subsection{FlashAttention forward}\label{sec:flash}

The same abstractions must cover the online-softmax FlashAttention forward
pass.  The sketch below is part of the spec; a conformance kernel implementing
it is required (Section~\ref{sec:conformance}).

\begin{lstlisting}
@tilegrad.jit
def flash_fwd(Q, K, V, *,
              Bq: T.Const[int] = 64,
              Bk: T.Const[int] = 64):
  B, Sq, Sk, D = T.const("B", "Sq", "Sk", "D")
  Q: T.Tensor[(B, Sq, D), T.float16]
  K: T.Tensor[(B, Sk, D), T.float16]
  V: T.Tensor[(B, Sk, D), T.float16]
  O  = T.empty((B, Sq, D), T.float16)
  Lse = T.empty((B, Sq), T.float32)        # log-sum-exp, for backward

  scale = (1.0 / D) ** 0.5

  with T.Kernel(B, T.ceildiv(Sq, Bq)) as (bb, qg):
    q_tile = T.shared((Bq, D), T.float16)
    k_tile = T.shared((Bk, D), T.float16)
    v_tile = T.shared((Bk, D), T.float16)
    s_frag = T.fragment((Bq, Bk), T.float32)
    acc    = T.fragment((Bq, D), T.float32)
    m_run  = T.fragment((Bq, 1), T.float32)   # running row-max
    l_run  = T.fragment((Bq, 1), T.float32)   # running row-sum
    o_run  = T.fragment((Bq, D), T.float32)   # running output

    T.clear(acc); T.clear(m_run); T.clear(l_run); T.clear(o_run)
    T.copy(Q.tile((bb, qg*Bq, 0), (Bq, D)), q_tile)

    for ko in T.Pipelined(T.ceildiv(Sk, Bk), stages=2):
      T.copy(K.tile((bb, ko*Bk, 0), (Bk, D)), k_tile)
      T.copy(V.tile((bb, ko*Bk, 0), (Bk, D)), v_tile)

      T.clear(s_frag)
      T.mma(q_tile, k_tile, s_frag)               # (Bq,D) @ (D,Bk)
      s_frag = s_frag * scale

      m_blk = T.reduce_max(s_frag, axis=-1, keepdim=True)
      m_new = T.maximum(m_run, m_blk)
      alpha = T.exp2(m_run - m_new)               # exp2 base for exp/log2
      beta  = T.exp2(s_frag - m_new)

      l_run = l_run * alpha + T.reduce_sum(beta, axis=-1, keepdim=True)
      o_run = o_run * alpha
      T.mma(beta, v_tile, o_run)                  # accumulate V contributions
      m_run = m_new

    o_run = o_run / l_run
    T.copy(o_run, O.tile((bb, qg*Bq, 0), (Bq, D)))
    T.copy(m_run + T.log2(l_run),
           Lse.tile((bb, qg*Bq), (Bq,)))
  return O, Lse
\end{lstlisting}

This example exercises every operation category \TileGrad{} commits to:
specialization symbols shared across multiple tensors, a 2D grid, managed
shared staging, collective fragments, elementwise fragment algebra, row
reductions with keepdim broadcasting, scalar-fragment mixing, two MMAs per
inner loop, a pipelined loop, and multi-output return.  If any of those is
missing or awkward, the spec has a hole.

\section{Program contract}

A decorated function consists of:

\begin{enumerate}
  \item runtime tensor and scalar parameters;
  \item keyword-only compile-time parameters (\texttt{T.Const});
  \item tensor and output declarations;
  \item exactly one \texttt{with T.Kernel(...)} launch region;
  \item zero or more compile-time helper or macro expansions; and
  \item a return of every declared output, preserving tuple order.
\end{enumerate}

\subsection{Value categories}

\begin{center}
\begin{tabularx}{\textwidth}{@{}lYY@{}}
\toprule
\textbf{Category} & \textbf{Meaning} & \textbf{Permitted uses} \\
\midrule
Compile-time constant & Included in the specialization cache key. &
Python control flow, tile and allocation shapes, workgroup shape, layouts,
unrolling, and all device expressions. \\
Tensor shape constant & Inferred from a runtime tensor, then specialized. &
Grid extents, output shapes, bounds, and device expressions. \\
Runtime scalar & Passed to the device program without specialization. &
Device expressions, predicates, indices, and loop bounds; not launch,
allocation, fragment, layout, unroll, or workgroup shapes. \\
\bottomrule
\end{tabularx}
\end{center}

Tensor rank and dtype are always specialization-time information.
\texttt{T.const(...)} declares named specialization symbols; each symbol must
occur as a dimension in at least one input tensor declaration, and repeated
occurrences must agree at specialization.  An unbound or conflicting symbol is
a specialization error.  Non-constant control flow and non-constant workgroup
shapes are rejected before device execution.

\subsection{Outputs}

\texttt{T.empty(shape, dtype)} declares a dense, contiguous global output
allocated by the host wrapper before launch.  Every output \must{} be returned
exactly once, in source order, as a tuple when there are several.  Outputs must
not alias one another or any input.  Runtime tensor inputs are read-only in V1;
an in-place extension would need to make aliasing explicit.  Scalar device
results are represented as rank-zero or one-element output tensors, not Python
values.

\subsection{Sizes, ranks, and forms}

Core implementations \must{} support tensor ranks zero through four and one- to
three-dimensional grids and workgroups.  External tensors are dense,
non-overlapping, and contiguous, matching the current \tinygrad{}
custom-kernel boundary.  Higher ranks and non-contiguous external tensors are
capabilities.

\section{Type and effect system}

\subsection{Memory spaces}

\begin{center}
\begin{tabularx}{\textwidth}{@{}lYYY@{}}
\toprule
\textbf{Space} & \textbf{Ownership} & \textbf{Shape} & \textbf{Purpose} \\
\midrule
Global & Invocation/tensor & Static rank; specialized dimensions &
Inputs and outputs. \\
Shared & Workgroup & Compile-time constant &
Cooperative staging and communication. \\
Fragment & Workgroup collective & Compile-time constant &
Logically shaped values distributed across lanes. \\
Private & One local agent & Compile-time constant &
Scalars and statically indexed register tuples. \\
\bottomrule
\end{tabularx}
\end{center}

Shared, fragment, and private allocations have lexical lifetime inside one
launch region.  Dynamic shared allocation is outside V1.  Private tuple
indices must become constant after specialization and unrolling; otherwise
compilation fails rather than silently spilling to another memory space.

\subsection{Views and bounds}

A buffer type is the tuple
\[
  (\mathrm{dtype},\;\mathrm{rank},\;\mathrm{shape},\;\mathrm{strides},
    \;\mathrm{address\ space},\;\mathrm{access}).
\]
A view is
\[
  V = (B,\;o,\;\mathbf{s},\;\mathbf{d},\;n_B),
\]
where \(B\) is a base buffer, \(o\) an element offset, \(\mathbf{s}\) a logical
shape, \(\mathbf{d}\) element strides, and the base storage contains \(n_B\)
elements.  Coordinate \(\mathbf{i}\) is in bounds exactly when
\[
  \mathrm{inbounds}(V,\mathbf{i}) \iff
  \left(\bigwedge_k 0\leq i_k<s_k\right)
  \land 0\leq\mathrm{addr}(V,\mathbf{i})<n_B,
  \qquad
  \mathrm{addr}(V,\mathbf{i}) = o + \textstyle\sum_k i_k d_k.
\]

\texttt{tile(origin, shape)} \may intentionally create a logical view extending
beyond base storage; masked load, store, fill, and copy use the predicate
above to define edge behavior.  \texttt{subview} and \texttt{transpose} are
view operations and do not copy.  \texttt{reshape} is valid only when
contiguity can be proven.  Writable zero-stride views, overlapping writable
views, and negative strides are invalid in V1.

Two views overlap when their sets of valid base-storage addresses intersect.
Copy between identical views is a no-op.  Every other overlapping source and
destination copy is invalid; V1 provides no snapshot or \texttt{memmove}
semantics, and compilation must fail when disjointness cannot be proven.

\subsection{Effects and synchronization}

Effects to the same logical location are ordered by source order unless they
occur in an unordered parallel region.  Two unordered writes, or an unordered
read and write without a defined collective or synchronization relation, form a
data race and are invalid; V1 has no atomic escape from this rule.

Managed high-level operations carry effect analysis and \must{} conservatively
preserve read-after-write, write-after-read before storage reuse,
write-after-write source order, and loop-carried dependencies.  A portable
workgroup barrier waits for all local agents and establishes a happens-before
relation from prior shared writes to subsequent shared reads by participants.
Every barrier \must{} be reached convergently by the same participant set.

Expert code \may{} request manually synchronized shared storage
(\texttt{sync="manual"}).  Managed and manual synchronization must not be mixed
for one allocation.  Raw shared access through explicit thread identifiers is
permitted only for manually synchronized allocations; the author is then
responsible for barriers and race freedom.  There is no grid-wide barrier in
V1.

\subsection{Uniformity}

A conditional is \emph{workgroup-uniform} when every local agent evaluates the
same branch.  Collective tile operations and barriers may appear only in
uniform control flow.  The compiler \must{} prove uniformity from compile-time
values, grid coordinates, and other known-uniform values, or reject the
program.  Local thread identifiers and values loaded under thread-varying
indices are conservatively non-uniform.  V1 has no user assertion that
overrides uniformity analysis.

\subsection{Scalar types}

Core scalar types are \texttt{bool}, signed and unsigned 8-, 16-, 32-, and
64-bit integers, \texttt{float16}, \texttt{bfloat16}, and \texttt{float32}.
Additional \tinygrad{} dtypes are capabilities.  Expressions are statically
typed; implicit conversion is limited to a Python literal converted to the type
required by its context.  All other dtype changes use \texttt{T.cast}.  Integer
indices use the implementation's index type and \mustnot{} be silently
converted from floating point.

\section{Operation categories}

This section names the operation categories V1 \must{} support.  Exact kwarg
signatures, defaults, and abbreviations live in the API reference; this section
defines what each category means and what the compiler guarantees about it.

\subsection{Loops}

\begin{center}
\begin{tabularx}{\textwidth}{@{}lY@{}}
\toprule
\textbf{Form} & \textbf{Semantics} \\
\midrule
Serial & Ordered iterations over a half-open integer interval. \\
Unroll & Serial semantics with a compile-time request to expand iterations. \\
Reduce & Ordered recurrence semantics that may reassociate under the
numerical contract. \\
Parallel & An unordered logical Cartesian iteration space; every point executes
exactly once. \\
Pipelined & Serial logical semantics with a positive compile-time pipeline
hint. \\
\bottomrule
\end{tabularx}
\end{center}

\subsection{Logical parallel mapping}

\texttt{T.Parallel} is not an alias for a hardware local axis.  It defines an
unordered logical iteration space distributed over the current workgroup.  Let
the iteration shape be \((e_0,\ldots,e_{r-1})\), let \(P\) be the workgroup
size, and let \(\ell\in[0,P)\) be the row-major linear local-agent identifier.
The compiler selects from a bounded set of canonical policies:

\begin{itemize}
  \item \texttt{row\_major}: \(q_R(\mathbf{i})=\sum_k i_k\prod_{j>k}e_j\);
        owner \(q_R\bmod P\), private slot \(\lfloor q_R/P\rfloor\);
  \item \texttt{column\_major}: \(q_C(\mathbf{i})=\sum_k i_k\prod_{j<k}e_j\);
        owner \(q_C\bmod P\), private slot \(\lfloor q_C/P\rfloor\); and
  \item an internal exact mapping required by a selected native matrix
        capability.
\end{itemize}

An agent executes the coordinates whose linearized value is \(\ell+tP\) for
nonnegative \(t\) while it is less than \(\prod_k e_k\).  The source may request
a canonical policy or \texttt{auto}; \texttt{auto} may select only from the
specified finite set, and selection \must{} appear in legalized debug IR.  If
explicit thread identifiers affect the same memory operations as a
\texttt{T.Parallel}, the loop policy \must{} be explicit; an implementation
\mustnot{} make observable behavior depend on an unreported mapping.

\subsection{Collective tile operations}

Tile operations are collective statements with explicit logical read and write
regions.  They \must{} be reached uniformly by every participating local agent;
a tile operation in thread-divergent control flow is invalid unless the
operation explicitly defines predicated participation.

\paragraph{Fill and clear.}
\(\op{Fill}(D,x)\) writes \(\mathrm{cast}_{D.\mathrm{dtype}}(x)\) to every
valid logical coordinate of \(D\).  \(\op{Clear}(D)=\op{Fill}(D,0)\).  Fill
never writes outside the destination view.

\paragraph{Synchronous copy.}
For equal-shaped logical regions \(S\) and \(D\),
\[
  \op{Copy}(S,D,p):\quad
  D[\mathbf{i}]\leftarrow
  \begin{cases}
    \mathrm{cast}_{D.\mathrm{dtype}}(S[\mathbf{i}]) & \mathrm{inbounds}(S,\mathbf{i})\\
    \mathrm{cast}_{D.\mathrm{dtype}}(p) & \text{otherwise}
  \end{cases}
\]
when \(D[\mathbf{i}]\) is in bounds.  Source and destination shapes must be
equal; copy does not broadcast and does not infer a tile extent from an
unrelated operand.  Copy carries safe-to-consume semantics: a subsequent
dependent high-level tile operation observes completed values.  Instruction
choice --- scalar, vector, SIMT-cooperative, or target-specific transfer --- is
not observable; a regular copy never exposes a partially completed asynchronous
transfer.

\paragraph{Matrix multiply-accumulate.}
For rank-two operands,
\[
  C \leftarrow C + \mathrm{cast}_{T_{acc}}(\mathrm{op}_A(A))
                    \times \mathrm{cast}_{T_{acc}}(\mathrm{op}_B(B)),
\]
with \(A:(M,K)\), \(B:(K,N)\), \(C:(M,N)\) after optional transpose, and no
implicit clear of \(C\).  Core implementations support \texttt{float32} SIMT
MMA and \texttt{float16}/\texttt{bfloat16} inputs with \texttt{float32}
accumulation; other dtype pairs are capabilities.  The selection mode is one of
\texttt{auto}, \texttt{simt}, or an exact native capability name.  An exact
request \mustnot{} silently fall back, and \texttt{auto} \must{} preserve the
same logical fragment and numerical contract across all selectable paths.

\paragraph{Axis reduction.}
\(\op{Reduce}(X,a,f,T_{acc},k)\) reduces axis \(a\) using operation \(f\)
(\texttt{sum}, \texttt{max}, or \texttt{min}), accumulator dtype \(T_{acc}\),
and keep-dimension flag \(k\).  Reduction starts from the operation identity
unless an explicit initial value is provided.  Maximum and minimum default to
\texttt{nan\_policy="propagate"}; an implementation may emulate this policy
when native instructions differ.  Reduction order is otherwise
implementation-defined.

\subsection{Fragment algebra}

A fragment is a statically shaped, workgroup-collective logical value; its
physical register ownership is not observable source state.  Access to a
fragment coordinate is legal only inside a tile operation or a matching
logical parallel region (equal extents and identical resolved policy).
Arbitrary fragment indexing by a runtime thread identifier is invalid.  SIMT
fallback \must{} implement the same logical coordinate ownership; it \mustnot{}
replicate the complete fragment in every thread.

\op{MapUnary}, \op{MapBinary}, \op{Select}, and \op{Cast} apply pointwise and
produce a fresh fragment.  Operands use right-aligned broadcasting: for each
aligned axis, extents must be equal or one; missing leading axes and
extent-one axes repeat the same source coordinate.  The result extent is the
maximum of the aligned extents.  Row and column broadcasting are not special
cases: shapes \((M,1)\) and \((1,N)\) follow the same rule.  Scalar operands
are rank-zero fragments.  All operations produce a new logical fragment
version; source versions are read before destination writes, including
syntactically in-place elementwise forms.

\subsection{Scalar memory}

Scalar load and store carry explicit masks and alternate values:
\begin{align*}
  \op{Load}(V,\mathbf{i},m,a) &=
    \begin{cases}V[\mathbf{i}] & m\land\mathrm{inbounds}(V,\mathbf{i})\\
                  a & \text{otherwise,}\end{cases}\\
  \op{Store}(V,\mathbf{i},x,m) &: \quad
    V[\mathbf{i}]\leftarrow x\quad\text{iff }m\land\mathrm{inbounds}(V,\mathbf{i}).
\end{align*}
Indexing abbreviations (\texttt{V[i]}) are unmasked forms of the same.
Floating-point reassociation and reduction order are implementation-defined
within the numerical contract; semantic equality is not generally bitwise
equality.  Detailed tolerance tables live in the conformance test suite.

\section{Capabilities}\label{sec:capabilities}

Capabilities describe optional legal lowerings, not alternate language
semantics.  They are the single mechanism by which a target advertises more
than scalar/SIMT fallback.  A capability record has:

\begin{itemize}
  \item a versioned, namespaced name;
  \item a target predicate;
  \item operand and accumulator dtypes;
  \item an atom shape \((M_a, N_a, K_a)\);
  \item a participant count;
  \item operand scope requirements (e.g., shared or private);
  \item the internal fragment layout the capability assumes;
  \item the exact \tinygrad{} \op{Wmma} argument schema used by the
        compatibility adapter.
\end{itemize}

A user requests a capability by name via the operation's mode argument, e.g.\
\texttt{T.mma(A, B, C, mode="tinygrad.cuda.m16n8k16.f16.f32.v1")}.  A backend
author publishes one via \texttt{tilegrad.capabilities(target)}, which returns
records of the following shape:

\begin{lstlisting}
Capability(
    name="tinygrad.cuda.m16n8k16.f16.f32.v1",
    target=("cuda", "sm>=80"),
    operand_dtypes=(float16, float16),
    accumulator_dtype=float32,
    atom=(16, 8, 16),
    participants=32,
    operand_scope=("shared", "shared"),
    layout="row_major_acc",
    wmma_schema=Ops.WMMA(arg_types=(...), ...),
)
\end{lstlisting}

A logical MMA \may use a native capability only when \(M\), \(N\), and \(K\) are
integer multiples of the corresponding atom dimensions.  The compiler
partitions \(M\) and \(N\) into disjoint atoms, visits \(K\) atoms in
increasing logical order, and suppresses out-of-bounds stores; the native MMA
itself never accesses a partial atom.  Renaming or changing the meaning of a
record requires a new versioned name; the adapter and record \must{} be jointly
tested against a pinned \tinygrad{} revision range.

The compiler \must{} distinguish an invalid program, a valid program
unavailable on the selected profile, a failure to generate valid \tinygrad{}
UOps, and a backend compilation or runtime failure.  Unsupported exact native
requests fail during capability legalization, before UOp construction, with a
diagnostic stating source construct, expected contract, observed types and
shapes, and selected target.

A V1 release \may{} publish zero required native capabilities; a conforming
core profile implementation already runs every program through SIMT fallback.
Conformance claims only require the capabilities the implementer advertises
under the names they advertise.

\section{Lowering contract}\label{sec:boundary}

Source programs pass through a small number of named stages:

\begin{center}
\begin{tabularx}{\textwidth}{@{}lY@{}}
\toprule
\textbf{Stage} & \textbf{Responsibility} \\
\midrule
Specialize & Bind compile-time constants and tensor shape constants. \\
Typed tile IR & Parse, type-check, verify shapes, effects, access modes,
uniformity, and resource rules. \\
Legalized tile IR & Select canonical logical-parallel and fragment ownership
layouts, plan synchronized managed memory, and select SIMT or exact native
capabilities. \\
\tinygrad{} UOps & Emit scalar, vector, range, memory, barrier, and optional
\op{Wmma} UOps; delegate rendering, compilation, launch, and storage. \\
\bottomrule
\end{tabularx}
\end{center}

Each named stage \must{} be inspectable at the corresponding named IR.  An
implementation \may{} add internal passes; their names are not part of the
contract.

The interface with \tinygrad{} is the UOp set \TileGrad{} \may{} emit.  V1
\must{} emit, and \tinygrad{} \must{} accept:

\begin{itemize}
  \item scalar ALU, compare, select, and cast nodes;
  \item \op{Load} and \op{Store} in global and local address spaces, with
        explicit mask support;
  \item \op{Loop}, \op{Range}, and reduction range nodes with their standard
        axis-type tags;
  \item thread and grid coordinate queries;
  \item \op{Barrier} with the workgroup memory scope; and
  \item optionally \op{Wmma}, schema-defined by an advertised capability.
\end{itemize}

\tinygrad{} \must{} preserve synchronization structure that \TileGrad{} emits:
a barrier \must{} remain a barrier with the same scope and participant set, and
\tinygrad{} \mustnot{} reorder a global \op{Store} across a barrier or eliminate
a declared output \op{Store}.  Within those constraints, \tinygrad{} \may{}
fuse, vectorize, lower to native instructions, schedule, and rewrite.  \TileGrad{}
\mustnot{} rely on \tinygrad{} optimization to recover synchronization, fragment
ownership, or memory dependences it did not emit.

\subsection{Differentiation}

V1 has no implicit autograd through a \TileGrad{} kernel.  Source-level
differentiation of a decorated kernel \must{} fail at JIT specialization with a
single clear diagnostic.  The supported differentiation paths are:

\begin{itemize}
  \item an explicit backward kernel written by the same author, returning VJP
        inputs in source order; and
  \item a registered VJP rule mapping a forward capability name to a backward
        capability name through the same capability mechanism.
\end{itemize}

V1 publishes the VJP registration interface, but \may{} ship with no registered
gradients.  The IR \must{} preserve what an autograd extension would need
(input views, output aliasing metadata, and capability names), so post-V1
differentiation is additive, not a rewrite.

\section{Done criteria}\label{sec:conformance}

A V1 release claims \emph{core conformance} by passing, on at least one
\tinygrad{} backend, the following programs:

\begin{enumerate}
  \item exact and masked one- through four-dimensional copy;
  \item internal strided subviews and transpose;
  \item source padding, destination suppression, cross-dtype copy, and rejected
        overlapping copies;
  \item one- through three-dimensional grids, canonical ownership mappings,
        logical parallel loops, and rejected runtime-scalar launch extents;
  \item uniform and rejected divergent collective operations;
  \item shared-memory producer/consumer and reuse dependencies;
  \item fragment unary, binary, broadcast, select, and cast operations;
  \item sum, maximum, and minimum reductions, including non-power-of-two tails;
  \item float32 SIMT GEMM with M, N, and K edge tiles;
  \item float16 and bfloat16 input GEMM with float32 accumulation;
  \item unavailable exact-native capability failures (negative test);
  \item multiple output allocation and return ordering; and
  \item a small causal and non-causal FlashAttention forward kernel using
        stable online softmax, edge sequence lengths, reduction, \texttt{exp2},
        and two MMA operations.
\end{enumerate}

The same mathematical cases \should{} run through SIMT and every advertised
native capability; negative tests are part of conformance, not optional
robustness.  An \emph{accelerated} profile is additive: it \may{} provide
vectorized copy, native matrix capabilities, real software pipelining, hidden
asynchronous instructions behind synchronous operations, and additional dtypes.
A backend \may{} implement some accelerated capabilities without implementing
all of them.

Numerical tolerances, including a dedicated FlashAttention tolerance row, are
normative inputs to the conformance test suite, not this document; they are
published alongside it and may be tightened in a patch release.  Integer,
boolean, movement, mask, and correctly-rounded cast tests are exact.

\section{Deferred, with reasons}

The following are deferred from V1.  Each is listed with a one-line rationale so
the deferral can be re-litigated without re-reading the whole spec.

\begin{itemize}
  \item Automatic schedule search and autotuning: schedule belongs to the
        author in V1; automatic search is a post-V1 layer over a stable IR.
  \item Arbitrary layout functions and swizzled layouts: V1 ships a bounded
        set of canonical logical-parallel and fragment layouts; swizzling
        arrives with a native capability that needs it.
  \item TMA, WGMMA, and TCGEN05: arrive only as named capabilities, never as
        source constructs.
  \item TileLang or TVM source compatibility: a different language; reuse is
        at the user-model level, not the AST level.
  \item Backend source injection and custom \tinygrad{} runtime backends:
        \tinygrad{} owns codegen and runtime.
  \item Clusters and tensor memory: cooperative multi-workgroup execution is
        not yet the right surface; arrive with a capability that defines them.
  \item User-managed asynchronous copy, wait groups, and warp specialization:
        V1 hides async behind the synchronous copy contract; exposing it
        directly would break the one-correctness-path rule.
  \item Scans and atomic memory operations: V1 rejects data races; atomics
        would need a defined memory-model section first.
  \item Implicit autograd: see Section~\ref{sec:boundary}.
  \item Grid-wide barriers: not exposed; programs that need them are invalid
        until a capability defines a legal collective grid reduction.
\end{itemize}

\appendix

\section{Normative IR nodes}

The exact in-memory Python classes are not normative.  A conforming typed tile
IR \must{} represent at least the following semantic nodes; their names are
referenced by the inspectable-stage contract in Section~\ref{sec:boundary}.

\begin{longtable}{@{}>{\raggedright\arraybackslash}p{3.0cm}>{\raggedright\arraybackslash}p{4.6cm}>{\raggedright\arraybackslash}p{8.0cm}@{}}
\toprule
\textbf{Node} & \textbf{Operands} & \textbf{Meaning} \\
\midrule
\endhead
Program & parameters, outputs, launch & One specialized callable program. \\
Buffer parameter & type, shape, access & Dense external tensor storage. \\
Scalar parameter & dtype, category & Compile-time or runtime scalar. \\
Output specification & shape, dtype, order & Host-allocated returned tensor. \\
View & base, offset, shape, strides, storage extent & Logical region sharing base storage. \\
Allocate shared & shape, dtype, sync mode & Workgroup-visible storage. \\
Allocate fragment & shape, dtype & Collective logical register tile. \\
Allocate private & shape, dtype & Per-agent statically indexed storage. \\
Kernel launch & grid, workgroup & Grid of independent workgroups. \\
Serial loop & start, stop, step, body & Ordered iteration. \\
Unroll loop & static interval, body & Ordered iteration requested to expand. \\
Reduce loop & interval, recurrence & Scalar recurrence. \\
Parallel loop & extents, policy, body & Unordered logical iteration. \\
Pipelined loop & interval, stages, body & Serial semantics with pipeline hint. \\
Thread index & dimension & Explicit local-agent coordinate or extent. \\
If & predicate, branches & Structured control flow. \\
Load & view, coordinate, mask, alternate & Predicated scalar read. \\
Store & view, coordinate, value, mask & Predicated scalar write. \\
Fill & destination, scalar & Collective typed fill. \\
Copy & source, destination, pad & Collective synchronous region copy. \\
Fragment map & operation, sources, broadcast shape & Pointwise collective fragment operation. \\
MMA & A, B, C, transpose, mode & Collective accumulating matrix multiply. \\
Reduce & source, axis, op, dtype, keepdim & Collective axis reduction. \\
Barrier & memory scope & Convergent workgroup synchronization. \\
\bottomrule
\end{longtable}

\end{document}